% Coprime Ramsey Numbers: Exact Values and Structural Properties
% E-JC submission version (formatted with e-jc.sty).
% Documentation reconciliation: 2026-10-17; see docs/coprime-ramsey-corrections.md.
%
% EJC papers *must* begin with the following two lines.
\documentclass[12pt]{article}
\usepackage[amsmath]{e-jc}

% amssymb, amsthm, hyperref are already loaded by e-jc.sty.
% Theorem-like environments (theorem, lemma, corollary, proposition,
% definition, conjecture, remark, open, problem, ...) are declared in
% e-jc.sty, so we do NOT redeclare them here.
% Do not change geometry/margins; the style file sets them.
\usepackage{mathtools}
\usepackage{booktabs}

\newcommand{\Rcop}{R_{\mathrm{cop}}}
\newcommand{\Pcop}{P_{\mathrm{cop}}}
\newcommand{\Ccop}{C_{\mathrm{cop}}}
\newcommand{\GRcop}{\mathrm{GR}_{\mathrm{cop}}}
\newcommand{\Gcop}{G_{\mathrm{cop}}}

% Public research draft: journal submission/publication has not been established.
% Override the journal template's placeholder dateline and citation footer.
\makeatletter
\def\the@dateline{Research draft, corrected October 17, 2026\\[0.5ex]\copy@right}
\def\ps@plain{%
  \def\@oddhead{}\def\@evenhead{}%
  \def\@oddfoot{\normalfont\hfil Research draft --- \thepage\hfil}%
  \let\@evenfoot\@oddfoot}
\makeatother
\pagestyle{plain}

% SUBJECT CODES (MSC 2020).
\MSC{05D10, 05C55, 11A05}

% COPYRIGHT: single author, CC BY.
\Copyright{The author. Released under the CC BY license (International 4.0).}

% TITLE (no footnotes on titles).
\title{Coprime Ramsey Numbers:\\
Exact Values and Structural Properties}

% AUTHOR.
\author{Alexander Richard Towell\authornote{1}}

\authortext{1}{Department of Computer Science, Southern Illinois University
  Edwardsville, Edwardsville, Illinois, USA
  (\email{lex@metafunctor.com}).
  ORCID: 0000-0001-6443-9897.}

\begin{document}

\maketitle

\begin{abstract}
We introduce the \emph{coprime Ramsey number} $\Rcop(k)$, defined as
the minimum $n$ such that every $2$-coloring of the edges of the coprime
graph on $\{1,\ldots,n\}$ contains a monochromatic copy of
$K_k$. Combining SAT-computed lower bounds with an elementary clique
upper bound, we establish $\Rcop(2)=2$, $\Rcop(3)=11$, and
$\Rcop(4)=59$, compute the $3$-color triangle number
$\Rcop(3;3)=53$, and determine exact values for path, cycle,
asymmetric, and Gallai variants. All clique-type values are prime.
This observation, first recorded here as a conjecture, has since been
proved in general by Du, Xi, Deng, and Ma, who showed
$\Rcop(k)=p_{R(k,k)-1}$, the $(R(k,k)-1)$-th prime, for the
edge-coloring variant studied in this paper. Their complete-pattern
transfer also covers Gallai triangles. Paths and cycles require separate
lower-bound arguments; the two longest path values remain historical
computational reports awaiting independent reproduction.
\end{abstract}

\smallskip\noindent\textbf{Keywords:} Ramsey theory, coprime graph,
SAT solving, graph coloring.

%% =================================================================
\section{Introduction}\label{sec:intro}
%% =================================================================

The \emph{coprime graph} $\Gcop(n)$ has vertex set
$[n]=\{1,2,\ldots,n\}$ with $\{i,j\}$ an edge whenever
$\gcd(i,j)=1$. This graph, studied by Erd\H{o}s and
S\'ark\"ozy~\cite{erdos1997coprime} and others~\cite{pomerance2015,
ahlswede1993}, encodes number-theoretic structure in a combinatorial
object: its edge density converges to $6/\pi^2\approx 0.608$, its
clique number equals $1+\pi(n)$ (where $\pi$ is the prime-counting
function), and its independence number is $\lfloor n/2\rfloor$ for
$n \ge 2$ (see Section~\ref{sec:structure} for short proofs).

We propose a Ramsey-theoretic study of this graph. Classical Ramsey
numbers $R(k,k)$~\cite{ramsey1930} concern the complete graph $K_n$;
replacing $K_n$ with
$\Gcop(n)$ yields a problem where Ramsey-type forcing must contend with
the arithmetic structure of coprimality.

The values reported here were computed and announced first (March
2026, in the repository accompanying this paper). Subsequently Du,
Xi, Deng, and Ma~\cite{du2026prime}, citing those computations, proved
a closed form for the clique-type numbers,
$\Rcop(k)=p_{R(k,k)-1}$; Remark~\ref{rem:duetal} details the
correspondence. Their Proposition~4.6 also covers Gallai triangles.
The path and cycle variants of Section~\ref{sec:variants} have
non-complete targets, for which the prime-clique upper bound remains
available but the matching lower-bound transfer can fail.

\begin{definition}\label{def:rcop}
The \emph{coprime Ramsey number} $\Rcop(k)$ is the minimum $n$ such
that every $2$-coloring of $E(\Gcop(n))$ contains a monochromatic
$K_k$. More generally, $\Rcop(k;c)$ is the minimum $n$ such that every
$c$-coloring of $E(\Gcop(n))$ contains a monochromatic $K_k$ (so
$\Rcop(k) = \Rcop(k;2)$).
\end{definition}

Since $E(\Gcop(n))\subseteq E(K_n)$, the restriction of any
$K_k$-avoiding coloring of $E(K_n)$ to the coprime edges is an
avoiding coloring of $\Gcop(n)$: every clique of $\Gcop(n)$ is a
clique of $K_n$ with the same edge colors. Taking $n = R(k,k)-1$
gives $\Rcop(k)\ge R(k,k)$; forcing
monochromatic cliques in the sparser coprime graph never requires
fewer vertices, and empirically requires many more for $k \ge 3$.
The same restriction argument shows that an avoiding coloring of
$\Gcop(n)$ restricts to one of $\Gcop(m)$ for every $m \le n$, so the
forcing property is upward closed in~$n$ and each $\Rcop$ value is
determined by the single SAT-to-UNSAT transition.

In the other direction, an elementary clique bound shows
$\Rcop(k)$ is well defined and pins it to within a prime:

\begin{proposition}[Clique upper bound]\label{prop:cone}
$\Rcop(k) \le p_{R(k,k)-1}$, where $p_m$ denotes the $m$-th prime.
More generally, for $c$ colors and targets $k_1,\ldots,k_c$, the
corresponding coprime Ramsey number is at most
$p_{R(k_1,\ldots,k_c)-1}$.
\end{proposition}

\begin{proof}
Let $n = p_{R(k,k)-1}$, so that $\pi(n) = R(k,k)-1$. The set
$\{1\} \cup \{p \le n : p \text{ prime}\}$ is pairwise coprime, hence
a clique of $\Gcop(n)$ of size $1 + \pi(n) = R(k,k)$. Any
$2$-coloring of $E(\Gcop(n))$ restricts to a $2$-coloring of the
edges of this $K_{R(k,k)}$, which contains a monochromatic $K_k$ by
the definition of $R(k,k)$. The multicolor and asymmetric cases are
identical.
\end{proof}

%% =================================================================
\section{Exact values}\label{sec:exact}
%% =================================================================

\begin{theorem}\label{thm:rcop2}
$\Rcop(2)=2$.
\end{theorem}

\begin{proof}
Since $\gcd(1,2)=1$, the graph $\Gcop(2)$ has a single edge, and any
$2$-coloring assigns it one color, producing a monochromatic $K_2$;
while $\Gcop(1)$ has no edges, so no monochromatic $K_2$ is forced at
$n=1$.
\end{proof}

\begin{theorem}\label{thm:rcop3}
$\Rcop(3)=11$.
\end{theorem}

\begin{proof}
By exhaustive incremental extension. Among all $2^{21}$ edge
$2$-colorings of $\Gcop(8)$ ($21$ coprime edges), exactly $36$ avoid
monochromatic triangles. We extend these through vertices $9$, $10$,
$11$: for each new vertex $v$, the coprime edges from $v$ to
$\{1,\ldots,v-1\}$ are enumerated, and every $2$-coloring of these new
edges is tested against each surviving avoiding coloring.

\begin{center}
\begin{tabular}{@{}cccc@{}}
\toprule
$n$ & Coprime edges & New edges at $n$ & Avoiding colorings \\
\midrule
8  & 21 & ---  & 36  \\
9  & 27 & 6  & 36  \\
10 & 31 & 4  & 156 \\
11 & 41 & 10 & 0   \\
\bottomrule
\end{tabular}
\end{center}

The increase from $36$ to $156$ at $n=10$ occurs because
$10=2\times 5$ is coprime to only $4$ predecessors, imposing weak
constraints. The collapse to $0$ at $n=11$ is driven by $11$ being
prime: it is coprime to all of $\{1,\ldots,10\}$, creating $10$ new
edges and far too many triangles for any $2$-coloring to avoid.
\end{proof}

\begin{remark}\label{rem:ratio3}
$\Rcop(3)/R(3,3)=11/6\approx 1.83$. The excess over $R(3,3)$
reflects multiplicative correlations among missing edges:
multiples of a prime $p$ form correlated non-adjacencies in
$\Gcop(n)$~\cite{erdos1997coprime}.
\end{remark}

\begin{theorem}\label{thm:rcop4}
$\Rcop(4)=59$.
\end{theorem}

\begin{proof}
Lower bound: for every $n\le 58$, the Glucose4
solver~\cite{glucose} finds a $2$-coloring of $E(\Gcop(n))$ with no
monochromatic $K_4$; by the restriction argument of
Section~\ref{sec:intro} this gives $\Rcop(4) \ge 59$.
Upper bound: $R(4,4) = 18$, and $59 = p_{17}$ is the $17$-th prime,
so Proposition~\ref{prop:cone} gives $\Rcop(4) \le p_{17} = 59$.
\end{proof}

\begin{remark}[Computational corroboration]\label{rem:rcop4-extension}
Prior to the observation of Proposition~\ref{prop:cone}, the upper
bound was explored using sampled extensions: $100$ distinct avoiding
colorings of $E(\Gcop(58))$ were extracted via SAT with blocking
clauses, and for each one the instance encoding the extension to
$n=59$ (coloring the $58$ new edges from the prime vertex~$59$) is
UNSAT. The full $n=59$ instance ($1{,}085$ variables,
${\sim}118{,}000$ clauses) did not terminate within $30$ minutes
under CDCL solving, which is why the seeded extensions were used;
these samples do not exhaust the avoiding colorings at $58$ and hence
do not prove the upper bound. Proposition~\ref{prop:cone} supplies the
complete upper-bound proof.
\end{remark}

\begin{center}
\begin{tabular}{@{}ccc@{}}
\toprule
$k$ & $\Rcop(k)$ & $R(k,k)$ \\
\midrule
2 & 2  & 2  \\
3 & 11 & 6  \\
4 & 59 & 18 \\
\bottomrule
\end{tabular}
\end{center}

\begin{theorem}\label{thm:rcop33}
$\Rcop(3;3)=53$.
\end{theorem}

\begin{proof}
A full $3$-coloring of $E(\Gcop(52))$ avoiding monochromatic
triangles gives the lower bound; a new witness was independently
checked against all $6{,}684$ coprime triangles in the October 17, 2026
audit. For the upper bound, $R(3,3,3)=17$ and
Proposition~\ref{prop:cone} give $\Rcop(3;3)\le p_{16}=53$.
The historical computation reported a direct UNSAT result at $53$;
that full search was not rerun in the audit and is not needed for
this upper-bound proof.
\end{proof}

\begin{conjecture}[Primality; now a theorem of Du--Xi--Deng--Ma~{\cite{du2026prime}}]\label{conj:prime}
All diagonal coprime Ramsey numbers $\Rcop(k)$ with $k \ge 2$ are
prime. More generally, $\Rcop(k;c)$ is prime for all $k, c \ge 2$.
\end{conjecture}

\noindent\emph{Evidence.} $\Rcop(2)=2$, $\Rcop(3)=11$, $\Rcop(4)=59$,
and $\Rcop(3;3)=53$ are all prime. The heuristic explanation is that
prime values of $n$ maximize the coprime density of $\Gcop(n)$---a prime
$p$ is coprime to all of $\{1,\ldots,p-1\}$---so the transition from
SAT to UNSAT preferentially occurs at primes, where the addition of a
universal coprime vertex creates the maximal number of new constraints.
This conjecture has since been proved by Du, Xi, Deng, and
Ma~\cite{du2026prime}; see Remark~\ref{rem:duetal} below.

\begin{remark}[Subsequent theoretical confirmation]\label{rem:duetal}
The exact values introduced and computed above have since been given a
closed-form explanation. Du, Xi, Deng, and Ma~\cite{du2026prime} study
coprime Ramsey numbers in two formulations, a vertex-coloring and an
edge-coloring version. The numbers $\Rcop(k)$ of the present paper,
defined by $2$-colorings of $E(\Gcop(n))$, are the \emph{edge}-coloring
version, and for it Du et al.\ prove (their Theorem~4.4) the prime-index
identity
\[
  \Rcop(k)=p_{R(k,k)-1},
\]
where $p_m$ denotes the $m$-th prime and $R(k,k)$ the classical diagonal
Ramsey number; the upper-bound half is exactly
Proposition~\ref{prop:cone}, and their contribution is the matching
lower-bound construction. Their vertex-coloring
version instead satisfies the elementary identity $p_{\sum_i(k_i-1)}$
(their Theorem~3.2), a different quantity. The values computed in this
paper are exactly the values the identity predicts: with $R(3,3)=6$
and $R(4,4)=18$ one obtains $p_{5}=11$ and $p_{17}=59$. Our
asymmetric and multicolor values are likewise consistent with the
mixed identity $p_{R(k_1,\ldots,k_c)-1}$ of their Theorem~4.4 in every
computed case: $\Rcop(2,3) = 3 = p_2$ with $R(2,3)=3$,
$\Rcop(2,4) = 5 = p_3$ with $R(2,4)=4$, $\Rcop(3,4) = 19 = p_8$ with
$R(3,4)=9$, and $\Rcop(3;3) = 53 = p_{16}$ with $R(3,3,3)=17$. The
primality observation of Conjecture~\ref{conj:prime} is thereby
established unconditionally: $\Rcop(k)$ is a prime by construction.
For the two-color diagonal, identifying the primes of the form
$p_{R(k,k)-1}$ remains tied to the unknown classical values. In the
unrestricted mixed family every prime occurs: $R(2,t)=t$ gives
$\Rcop(2,t)=p_{t-1}$. Their Proposition~4.6 also transfers fixed color patterns on
complete witnesses. In particular, monochromatic or rainbow triangles
give $\GRcop(3;3)=p_{\mathrm{gr}_3(K_3)-1}=p_{10}=29$, since
$\mathrm{gr}_3(K_3)=11$. Paths and cycles have non-complete targets;
labels can coincide on non-adjacent vertices, so the matching pullback
lower bound is unavailable in general. Their non-prime values
($\Pcop$ attains $9$ and $10$, and $\Ccop$ attains $8$) demonstrate
this boundary.
\end{remark}

%% =================================================================
\section{Ramsey variants on the coprime graph}\label{sec:variants}
%% =================================================================

We define and compute several Ramsey-type variants on $\Gcop(n)$. All
values are historical SAT computations. The October 17, 2026 audit
reproduced the smaller exact transitions, but the path cases $k=7,8$
each exceeded a $60$-second audit limit; those two reported values
remain awaiting independent reproduction. Clique and Gallai triangle
values also follow from the complete-pattern transfer.

\begin{definition}\label{def:variants}
Let $c\ge 2$ be the number of colors and $k\ge 2$ the target structure
size.
\begin{enumerate}
\item $\Pcop(k)$: minimum $n$ such that every $2$-coloring of
  $E(\Gcop(n))$ contains a monochromatic path on $k$ edges.
\item $\Ccop(k)$: minimum $n$ such that every $2$-coloring of
  $E(\Gcop(n))$ contains a monochromatic cycle on $k$ edges.
\item $\Rcop(s,t)$: minimum $n$ such that every $2$-coloring of
  $E(\Gcop(n))$ has a monochromatic $K_s$ in color~$0$ or $K_t$ in
  color~$1$ (asymmetric coprime Ramsey).
\item $\GRcop(k;c)$: the minimum $n$ such that every $c$-coloring of
  $E(\Gcop(n))$ contains a monochromatic or rainbow $K_k$ (the Gallai
  coprime Ramsey number; requires $c\ge 3$).
\end{enumerate}
\end{definition}

\begin{table}[ht]
\centering
\caption{Reported coprime Ramsey variants. The path values for $k=7,8$
await independent reproduction in the current audit.}\label{tab:variants}
\medskip
\begin{tabular}{@{}lll@{}}
\toprule
\textbf{Variant} & \textbf{Values} & \textbf{Classical counterpart} \\
\midrule
$\Pcop(k)$, $k=3,\ldots,8$ & $5,\; 7,\; 9,\; 10,\; 13,\; 13$ & --- \\[3pt]
$\Ccop(k)$, $k=3,\ldots,6$ & $11,\; 8,\; 13,\; 11$ & --- \\[3pt]
$\Rcop(2,3)$ & $3$ & $R(2,3) = 3$ \\
$\Rcop(2,4)$ & $5$ & $R(2,4) = 4$ \\
$\Rcop(3,4)$ & $19$ & $R(3,4) = 9$ \\[3pt]
$\Rcop(3;3)$ & $53$ & $R(3,3,3) = 17$ \\[3pt]
$\GRcop(3;3)$ & $29$ & $\mathrm{gr}_3(K_3) = 11$ \\
\bottomrule
\end{tabular}
\end{table}

\begin{remark}\label{rem:observations}
Several patterns emerge from Table~\ref{tab:variants}.
\begin{enumerate}
\item \emph{Path versus clique.} $\Pcop(k)\ll\Rcop(k)$: paths are much
  easier to force in the computed cases. The reported values plateau
  at $13$ for $k=7,8$; that plateau has not yet been independently
  reproduced by the current audit.
\item \emph{Even--odd cycle gap.} $\Ccop(4)=8<\Ccop(3)=11$ and
  $\Ccop(6)=11<\Ccop(5)=13$: even-length cycles are easier to force.
\item \emph{Multi-color growth.}
  $\Rcop(3;3)/\Rcop(3;2)=53/11\approx 4.8$, compared with the
  classical $R(3,3,3)/R(3,3)=17/6\approx 2.8$. The coprime graph's
  structure amplifies multi-color avoidance.
\item \emph{Coprime amplification.} For the clique-type variants
  with $k \ge 3$, the coprime version strictly exceeds its classical
  counterpart: $\Rcop(3)>R(3,3)$, $\Rcop(3;3)>R(3,3,3)$,
  $\GRcop(3;3)>\mathrm{gr}_3(K_3)$. (For $k = 2$ equality holds:
  $\Rcop(2) = 2 = R(2,2)$.)
\end{enumerate}
\end{remark}

%% =================================================================
\section{Structural properties}\label{sec:structure}
%% =================================================================

For completeness we record short proofs of the structural parameters
quoted in the introduction.

\begin{proposition}\label{prop:clique}
$\omega(\Gcop(n))=1+\pi(n)$.
\end{proposition}

\begin{proof}
The set $\{1\}\cup\{p\le n:p\text{ prime}\}$ is a clique of size
$1+\pi(n)$. For the upper bound: in any pairwise coprime set
$S\subseteq\{2,\ldots,n\}$, each element's prime factorization uses a
disjoint set of primes, so $|S|\le\pi(n)$.
\end{proof}

\begin{proposition}\label{prop:perfect}
$\Gcop(n)$ is weakly perfect: $\chi(\Gcop(n))=\omega(\Gcop(n))
=1+\pi(n)$.
\end{proposition}

This follows from the observation that the partition
\[
  \{1\},\ \{2,4,6,\ldots\},\ \{3,9,15,\ldots\},\ \{5,25,35,\ldots\},\ \ldots
\]
into classes indexed by the smallest prime factor gives a proper
$(1+\pi(n))$-coloring. Vertex colorings and chromatic numbers of
related coprime-type graphs on group elements were studied by
Syarifudin and Wardhana~\cite{syarifudin2021}.

\begin{proposition}\label{prop:independence}
For $n \ge 2$, $\alpha(\Gcop(n))=\lfloor n/2\rfloor$.
\end{proposition}

\begin{proof}
The even numbers $\{2,4,6,\ldots,2\lfloor n/2\rfloor\}$ form an
independent set of size $\lfloor n/2\rfloor$ (every pair shares the
factor~$2$). Conversely, let $S$ be an independent set with
$|S| \ge \lfloor n/2\rfloor+1 \ge 2$. Then $1 \notin S$, since $1$ is
adjacent to every other vertex, and $S$ contains no two consecutive
integers, since consecutive integers are coprime. But a subset of
$\{2,\ldots,n\}$ with no two consecutive elements has size at most
$\lceil (n-1)/2 \rceil = \lfloor n/2 \rfloor$, a contradiction.
\end{proof}


%% =================================================================
\section{Methodology}\label{sec:method}
%% =================================================================

All computations use the PySAT toolkit~\cite{pysat} with the
Glucose4~\cite{glucose} CDCL solver. The methodology varies by problem
size.

\paragraph{Small two-color instances ($\Rcop(3)$, path and cycle variants).}
For each $n$, we encode the avoidance problem as a SAT instance: one
Boolean variable per coprime edge ($x_e=\top$ for color~$0$, $x_e=\bot$
for color~$1$), two clauses per target substructure (forbidding
all-color-$0$ and all-color-$1$). The solver returns SAT or UNSAT. The
value $\Rcop(3)=11$ is also verified by the exact incremental extension
method of Theorem~\ref{thm:rcop3}, which enumerates \emph{all} avoiding
colorings at each $n$ and is thus a complete proof.

\paragraph{Large instances ($\Rcop(4)$).}
The incremental SAT approach determines $\Rcop(4)\ge 59$: Glucose4
returns SAT for all $n\le 58$, with solving times under $1$~second per
instance. The upper bound is Proposition~\ref{prop:cone}; the seeded
extension experiment of Remark~\ref{rem:rcop4-extension} provides
additional corroboration.

\paragraph{Multi-color variants ($\Rcop(3;3)$, $\GRcop(3;3)$).}
For $c$ colors, each edge has $c-1$ Boolean variables with
at-most-one pairwise exclusion constraints. For each clique and each
color, a clause forbids monochromatic assignment. The Gallai variant
additionally forbids rainbow cliques via all-distinct permutation
clauses. These instances are within Glucose4's reach for the parameter
ranges in Table~\ref{tab:variants}.

%% =================================================================
\subsection*{Acknowledgements}
%% =================================================================

Computational analysis was performed with the assistance of an AI
system (Claude, Anthropic), used for exploring the problem space and
generating candidate SAT encodings. Codex (OpenAI) assisted the
October 17, 2026 audit and subsequent software and documentation repairs.
The audit
independently checked full lower-bound witnesses at $10$, $58$, and $52$
and reproduced the smaller variant transitions; the two longest path
reports remain qualified as described in Section~\ref{sec:variants}.
No Lean certification of the coprime Ramsey results is claimed.
The author takes full responsibility for the correctness of the
statements and proofs. The author thanks Terry Tao and the contributors
to the Erd\H{o}s problems database~\cite{tao2025erdos} for creating a
resource that motivated this investigation.

\smallskip\noindent
Code and data are available at
\href{https://github.com/queelius/computational-explorations}%
{\nolinkurl{github.com/queelius/computational-explorations}}.


%% =================================================================
\begin{thebibliography}{12}

\bibitem{ahlswede1993}
R.~Ahlswede and L.~H.~Khachatrian.
Classical results on primitive and recent results on cross-primitive
sequences.
In \emph{The Mathematics of Paul Erd\H{o}s I}, pp.~104--116.
Springer, 1997.

\bibitem{du2026prime}
Z.~Du, W.~Xi, Z.~Deng, and L.~Ma.
Prime certificates for exact vertex-coprime Ramsey numbers.
\arxiv{2605.26815}, 2026.

\bibitem{erdos1997coprime}
P.~Erd\H{o}s and G.~N.~S\'ark\"ozy.
On cycles in the coprime graph of integers.
\emph{Electron.\ J.\ Combin.}, 4(2):R8, 1997.

\bibitem{glucose}
G.~Audemard and L.~Simon.
Predicting learnt clauses quality in modern {SAT} solvers.
In \emph{Proc.\ IJCAI}, pp.~399--404, 2009.

\bibitem{pysat}
A.~Ignatiev, A.~Morgado, and J.~Marques-Silva.
{PySAT}: A {P}ython toolkit for prototyping with {SAT} oracles.
In \emph{Proc.\ SAT}, pp.~428--437. Springer, 2018.

\bibitem{pomerance2015}
C.~Pomerance and J.~L.~Selfridge.
Proof of {D.~J.~Newman's} coprime mapping conjecture.
\emph{Mathematika}, 27:69--83, 1980.

\bibitem{ramsey1930}
F.~P.~Ramsey.
On a problem of formal logic.
\emph{Proc.\ London Math.\ Soc.\ (2)}, 30:264--286, 1930.

\bibitem{syarifudin2021}
A.~G.~Syarifudin and I.~G.~A.~W.~Wardhana.
The vertex coloring and chromatic number of coprime graph of the
dihedral group.
\emph{J.\ Phys.: Conf.\ Ser.}, 1783:012006, 2021.

\bibitem{tao2025erdos}
T.~Tao et~al.
Erd\H{o}s problems database.
\href{https://github.com/teorth/erdosproblems}%
{\nolinkurl{github.com/teorth/erdosproblems}}, 2025.

\end{thebibliography}

\end{document}
